\documentclass{article}

\title{Digital Systems Lab \\ \textlatin{\textbf{Project 1 \\ 7-Segment Display Driver}}}
\date{November 2026}

\author{
  \IfFileExists{UTH_greek_logo.png}{\includegraphics[width=0.2\linewidth]{UTH_greek_logo.png} \\ \vspace*{0.3in}}{}%
  \begin{tabular}{c}\textbf{Tsiantos Dimitrios} \\
  \texttt{\textlatin{dtsiantos@uth.gr}} \\
    3796 \\
  \end{tabular} \\
}

\usepackage[letterpaper, top=2cm, bottom=2cm, left=3cm, right=3cm, marginparwidth=1.75cm]{geometry}
\usepackage{graphicx}
\usepackage{csquotes}
\usepackage{listings}
\usepackage{xcolor}
\usepackage{tabularx}
\usepackage{tikz}
\usepackage{colortbl}
\usepackage{mathtools}
\usepackage[table]{xcolor}
\usepackage{hyperref}
\usepackage{amsmath}
\usepackage{subcaption}
\usepackage{algorithm}
\usepackage{algpseudocode}
\usepackage{booktabs}
\usepackage{siunitx}
\usepackage{rotating}
\usepackage{array}
\usepackage{makecell}
\usepackage{placeins}
\usepackage{url}

\usepackage{babel}
\babelprovide[import]{greek}
\babelprovide[import]{english}
\usepackage{fontspec}
\setmainfont{Libertinus Serif}
\setsansfont{Libertinus Sans}

% CODE LISTINGS
\definecolor{codegreen}{rgb}{0,0.6,0}
\definecolor{codegray}{rgb}{0.5,0.5,0.5}
\definecolor{codepurple}{rgb}{0.58,0,0.82}
\definecolor{backcolour}{rgb}{0.95,0.95,0.92}

\lstdefinestyle{mystyle}{
    backgroundcolor=\color{backcolour},   
    commentstyle=\color{codegreen},
    keywordstyle=\color{magenta},
    numberstyle=\tiny\color{codegray},
    stringstyle=\color{codepurple},
    basicstyle=\ttfamily\footnotesize,
    breakatwhitespace=false,         
    breaklines=true,                 
    captionpos=b,                    
    keepspaces=true,                 
    numbers=left,                    
    numbersep=5pt,                  
    showspaces=false,                
    showstringspaces=false,
    showtabs=false,                  
    tabsize=2
}

\lstset{style=mystyle}
% END CODE LISTINGS

\begin{document}

\maketitle
\clearpage

\tableofcontents
\clearpage

\section{Part A: Decoder}
\subsection{The module}
The \texttt{LEDdecoder} module receives an input of four bits named \texttt{char}
and outputs a signal \texttt{LED} of seven bits, each corresponding to a cathode
of a 7-Segment Display.

This module is combinational, meaning it does not use
any memory elements and the output only depends on the current input. This
can be implemented in Verilog using a combinational \texttt{always} block; that
is an \texttt{always} block whose sensitivity list doesn't contain the clock signal.

Depending of the value of \texttt{char}, the \texttt{case} statement combinationally
assigns the seven bit constant of the right hand side to the \texttt{LED} signal.

\lstinputlisting[language=Verilog]{LEDdecoder.v}

\subsection{Include directive}
The \texttt{`include} directive inserts the contents of the given file in place,
before compilation~\cite{ieee1364}. Vivado searches for included files in the
directory of the including source and in the include search paths of the
project~\cite{ug901}. The constants are therefore defined as \texttt{`define}
macros with an include guard, so that the header can be included by several
modules without redefinition.

\lstinputlisting[language=Verilog]{constants.vh}

\subsection{Case statement}
Since we have an input of four bits, we can only index $2^4 = 16$ different characters.
The \texttt{default} case should never be reached since all cases are covered
~\cite[Case Statements]{ug901}, but it is a good idea to use a default in the event
some case is accidentally forgotten later.

\subsection{Sensitivity list}
The sensitivity list only contains the input \texttt{char}, which, in this case, is equivalent
to using the "star" operator (\texttt{always @(*)}).

\subsection{Timescale}
The \texttt{`timescale} directive has the form
\texttt{`timescale} \texttt{<time\_unit>}\texttt{/}\texttt{<time\_precision>}, in this case \texttt{1ns/1ps}.
The time unit is the unit of all delays in the module (so \texttt{\#5} means
\SI{5}{\nano\second}), while the time precision is the smallest step to which
delays are rounded. It only affects simulation and is ignored by synthesis.
The Vivado simulator runs with the finest precision specified among all source
files of the design~\cite{ug900}. The same timescale is used in all verilog files
in this project.

\section{Part B: Four Digits Driver}
\subsection{MMCM Instantiation}

The board provides a \SI{100}{\mega\hertz} clock (period \SI{10}{\nano\second}),
while the anode and segment signals must stay stable for much longer, because
of the large capacitance of the LED circuit. Instead of dividing the clock with
a counter in the fabric, we use a Mixed-Mode Clock Manager (MMCM), a dedicated
clock management block of the 7~Series devices~\cite{ug472}. The MMCM produces
a new, slower clock with a fixed relationship to the input clock. All the sequential
logic of the driver then runs on this clock, and its outputs change once per clock cycle.

We instantiate the \texttt{MMCME2\_BASE} primitive and use the following
ports~\cite{ug472}:
\begin{description}
  \item[\texttt{CLKIN1}] The input clock, here the \SI{100}{\mega\hertz} board clock.
  \item[\texttt{CLKFBOUT}, \texttt{CLKFBIN}] The feedback loop. \texttt{CLKFBOUT}
        is connected to \texttt{CLKFBIN}, so that the MMCM can compare the
        feedback clock with the input clock and lock its internal oscillator (VCO).
  \item[\texttt{CLKOUT1}] The output clock used by our design.
  \item[\texttt{LOCKED}] Goes high when the output clocks are stable. Until then
        the design is held in reset (\texttt{rst\_int = reset\_db | \textasciitilde locked}).
  \item[\texttt{RST}] Asynchronous reset of the MMCM, driven by the debounced reset button.
  \item[\texttt{PWRDWN}] Powers down the MMCM; tied to 0.
\end{description}
The remaining outputs (\texttt{CLKOUT0}, \texttt{CLKOUT2}--\texttt{CLKOUT6}
and the inverted outputs) are left unconnected.

The MMCM first multiplies the input frequency to drive its VCO and then divides
the VCO frequency separately for each output~\cite{ug472}:
\begin{equation}
  f_{VCO} = f_{IN} \cdot \frac{M}{D}, \qquad
  f_{OUT} = \frac{f_{VCO}}{O}
\end{equation}
where $M$ = \texttt{CLKFBOUT\_MULT\_F}, $D$ = \texttt{DIVCLK\_DIVIDE} and
$O$ = \texttt{CLKOUT1\_DIVIDE}. The VCO frequency must lie within the range of
the device (\SIrange{600}{1200}{\mega\hertz} for speed grade $-1$~\cite{ds189}),
so the clock cannot be divided by 20 directly. We choose $M = 6$, $D = 1$ and
$O = 120$:
\begin{equation}
  f_{VCO} = \SI{100}{\mega\hertz} \cdot \frac{6}{1} = \SI{600}{\mega\hertz},
  \qquad
  f_{OUT} = \frac{\SI{600}{\mega\hertz}}{120} = \SI{5}{\mega\hertz}
  \;\Rightarrow\; T_{OUT} = \SI{200}{\nano\second}.
\end{equation}
\texttt{CLKIN1\_PERIOD} is set to \SI{10.0}{\nano\second} to inform the tools
of the input clock period.


\section{Part C: Message Rotation Using a Button}
\section{Part D: Message Rotation Using a Delay}

\bibliographystyle{IEEEtran}
\bibliography{references}

\end{document}
